% Part: incompleteness
% Chapter: representability-in-q
% Section: representing-relations

\documentclass[../../../include/open-logic-section]{subfiles}

\begin{document}

\olfileid{inc}{req}{rel}

\olsection{સંબંધોનું નિરૂપણ}

\emph{સંબંધ} નિરૂપ્ય હોય તેનો
અર્થ શું છે તે કહીએ.

\begin{defn}
\ollabel{defn:representing-relations}
પ્રાકૃતિક સંખ્યાઓ પરનો સંબંધ
$R(x_0,\dots,x_k)$ {\em
$\Th{Q}$માં નિરૂપ્ય} કહેવાય,
જો એવું સૂત્ર $!A_R(x_0,\dots,x_k)$
હોય કે $R(n_0,\dots,n_k)$
સાચું હોય ત્યારે
$\Th{Q}$માં
$!A_R(\num{n_0},\dots,\num{n_k})$
સાબિત થાય, અને
$R(n_0,\dots,n_k)$ ખોટું હોય
ત્યારે $\Th{Q}$માં
$\lnot !A_R(\num{n_0},
\dots, \num{n_k})$ સાબિત થાય.
\end{defn}

\begin{thm}
\ollabel{thm:representing-rels}
સંબંધ $\Th{Q}$માં નિરૂપ્ય
હોય ત્યારે અને માત્ર ત્યારે જ
તે સંગણનીય હોય છે.
\end{thm}

\begin{proof}
પહેલી દિશા માટે માનો કે
$R(x_0,\dots,x_k)$નું નિરૂપણ
$!A_R(x_0,\dots,x_k)$ સૂત્ર
કરે છે. $R$ ગણવા માટેની
રીત આ છે: ઇનપુટ
$n_0$, \dots,~$n_k$ મળે ત્યારે
$!A_R(\num{n_0}, \dots, \num{n_k})$
ની સાબિતી અને
$\lnot !A_R(\num{n_0}, \dots, \num{n_k})$
ની સાબિતી એકસાથે શોધો.
આપણી ધારણા પ્રમાણે એ બેમાંથી
એક સાબિતી જરૂર મળશે; પહેલી
મળે તો ``હા'' કહો, અન્યથા
``ના'' કહો.

બીજી દિશામાં, માનો કે
$R(x_0, \dots, x_k)$ સંગણનીય
છે. વ્યાખ્યા પ્રમાણે આનો અર્થ
તેનું લાક્ષણિક વિધેય
$\Char{R}(x_0, \dots, x_k)$
સંગણનીય છે. \olref[int]{thm:representable-iff-comp}
પ્રમાણે, $\Char{R}$ને કોઈ
સૂત્ર, કહો કે
$!A_{\Char{R}}(x_0, \dots, x_k,
y)$, નિરૂપે છે.
$!A_R(x_0, \dots, x_k)$ તરીકે
$!A_{\Char{R}}(x_0,
\dots, x_k, \num{1})$ સૂત્ર
લો. હવે કોઈપણ $n_0$,
\dots,~$n_k$ માટે,
$R(n_0, \dots, n_k)$ સાચું
હોય તો $\Char{R}(n_0, \dots,
n_k) = 1$; આ કિસ્સામાં
$\Th{Q}$માં
$!A_{\Char{R}}(\num{n_0}, \dots, \num{n_k},
\num{1})$ સાબિત થાય છે,
એટલે $\Th{Q}$માં
$!A_R(\num{n_0}, \dots,
\num{n_k})$ સાબિત થાય છે.
બીજી તરફ, $R(n_0, \dots,
n_k)$ ખોટું હોય તો
$\Char{R}(n_0, \dots, n_k) = 0$.
આથી $\Th{Q}$માં આ સાબિત થાય છે:
\[
\lforall[y][(!A_{\Char{R}}(\num{n_0}, \dots, \num{n_k}, y) \lif y =
  \num{0})].
\]
$\Th{Q}$માં
$\eq/[\num{0}][\num{1}]$
સાબિત થાય છે, તેથી
$\Th{Q}$માં
$\lnot !A_{\Char{R}}(\num{n_0}, \dots, \num{n_k}, \num{1})$
સાબિત થાય છે; પરિણામે
$\lnot !A_R(\num{n_0}, \dots, \num{n_k})$
પણ સાબિત થાય છે.
\end{proof}

\begin{prob}
$R$ એ~$\Th{Q}$માં નિરૂપ્ય
હોય, તો~$\Char{R}$ પણ તેમાં
નિરૂપ્ય છે તે બતાવો.
\end{prob}

\end{document}
